%%% BEGIN UNIT
\wbunit{Fracciones que valen lo mismo}{Vas a descubrir que dos fracciones diferentes pueden valer lo mismo y cómo encontrarlas.}

En las unidades anteriores aprendiste que una fracción nombra partes iguales de un entero. Ahora vas a comparar fracciones y verás que algunas se escriben diferente, pero representan la misma cantidad.

\section{Descubrir fracciones equivalentes con tiras y círculos}

Imagina dos tiras de papel del mismo tamaño. En la primera coloreas $\frac{1}{2}$. En la segunda divides cada mitad en 2 partes iguales y coloreas 2 de 4 partes. Las dos tiras tienen la misma cantidad coloreada.

\begin{center}
\wbfracbar[6cm]{1}{2}\\[2pt]
\wbfracbar[6cm]{2}{4}
\end{center}

La fracción cambió de nombre, pero no de cantidad:
\[
\frac{1}{2}=\frac{2}{4}
\]

Con los círculos pasa lo mismo. Un tercio y dos sextos cubren la misma parte de un círculo del mismo tamaño.

\begin{center}
\wbfraccircle[1.2cm]{1}{3}\hspace{1cm}\wbfraccircle[1.2cm]{2}{6}
\end{center}

\[
\frac{1}{3}=\frac{2}{6}
\]

También puedes verlo con tiras de tres y de seis partes.

\begin{center}
\wbfracbar[6cm]{2}{3}\\[2pt]
\wbfracbar[6cm]{4}{6}
\end{center}

\[
\frac{2}{3}=\frac{4}{6}
\]

\begin{wbidea}[Fracciones equivalentes]
Dos fracciones son equivalentes cuando representan la misma cantidad de un entero del mismo tamaño. Se escriben con el signo igual:
\[
\frac{1}{2}=\frac{2}{4}
\]
\end{wbidea}

\begin{wbexample}[¿Un medio y dos cuartos valen lo mismo?]
\textbf{Paso 1.} Dibuja dos tiras del mismo tamaño.

\textbf{Paso 2.} Divide la primera tira en 2 partes iguales y colorea 1. Eso es $\frac{1}{2}$.

\textbf{Paso 3.} Divide la segunda tira en 4 partes iguales y colorea 2. Eso es $\frac{2}{4}$.

\textbf{Paso 4.} Compara las partes coloreadas. Cubren la misma cantidad.

\begin{center}
\wbfracbar[6cm]{1}{2}\\[2pt]
\wbfracbar[6cm]{2}{4}
\end{center}

\textbf{Respuesta:} Sí, $\frac{1}{2}=\frac{2}{4}$.
\end{wbexample}

\begin{wbexample}[Un tercio y dos sextos]
\textbf{Paso 1.} Dibuja dos círculos del mismo tamaño.

\textbf{Paso 2.} Divide el primer círculo en 3 partes iguales y colorea 1. Eso es $\frac{1}{3}$.

\textbf{Paso 3.} Divide el segundo círculo en 6 partes iguales y colorea 2. Eso es $\frac{2}{6}$.

\textbf{Paso 4.} Observa la parte coloreada. Es la misma cantidad.

\begin{center}
\wbfraccircle[1.2cm]{1}{3}\hspace{1cm}\wbfraccircle[1.2cm]{2}{6}
\end{center}

\textbf{Respuesta:} $\frac{1}{3}=\frac{2}{6}$.
\end{wbexample}

\begin{wbtip}[Cuidado]
Las fracciones equivalentes deben ser parte de enteros del mismo tamaño. $\frac{1}{2}$ de una pizza grande no es equivalente a $\frac{1}{2}$ de una pizza pequeña. Además, $\frac{2}{4}$ no es más grande que $\frac{1}{2}$; solo tiene más partes.
\end{wbtip}

\begin{wbpractice}
\begin{enumerate}
\item Observa las tiras. ¿Son equivalentes? \wblines{1}

\begin{center}
\wbfracbar[5cm]{1}{2}\\[2pt]
\wbfracbar[5cm]{2}{4}
\end{center}

\item Colorea $\frac{1}{3}$ en el primer círculo y $\frac{2}{6}$ en el segundo. ¿Cubren la misma cantidad? \wblines{2}

\begin{center}
\wbfraccircle[1.2cm]{0}{3}\hspace{1cm}\wbfraccircle[1.2cm]{0}{6}
\end{center}

\item Completa: $\frac{1}{2}=\frac{\wbblank[1cm]}{4}$. \wblines{1}

\item Completa: $\frac{1}{3}=\frac{\wbblank[1cm]}{6}$. \wblines{1}

\item Compara: $\frac{1}{2}$ \wbblank[1.5cm] $\frac{2}{4}$. \wblines{1}

\item Ana pinta $\frac{1}{2}$ de una cartulina y Luis pinta $\frac{2}{4}$ de otra cartulina del mismo tamaño. ¿Quién pintó más? Explica. \wblines{3}

\item Explica con tus palabras qué significa que dos fracciones sean equivalentes. \wblines{3}
\end{enumerate}
\end{wbpractice}

\section{Doblar y partir para encontrar equivalentes}

Puedes encontrar fracciones equivalentes si divides cada parte en el mismo número de partes pequeñas.

Si divides cada mitad en 2 partes iguales, obtienes 4 partes en total y 2 partes coloreadas:
\[
\frac{1}{2}=\frac{2}{4}
\]

Si divides cada tercio en 2 partes iguales, obtienes 6 partes en total y 2 partes coloreadas:
\[
\frac{1}{3}=\frac{2}{6}
\]

\begin{center}
\wbfracbar[6cm]{1}{2}\\[2pt]
\wbfracbar[6cm]{2}{4}\\[6pt]
\wbfracbar[6cm]{1}{3}\\[2pt]
\wbfracbar[6cm]{2}{6}
\end{center}

Con rectángulos también se ve. Dos cuartos y cuatro octavos cubren la misma cantidad.

\begin{center}
\wbfracgrid{2}{1}{4}\hspace{0.8cm}\wbfracgrid{4}{2}{4}
\end{center}

\[
\frac{2}{4}=\frac{4}{8}
\]

\begin{wbidea}[Partir cada parte]
Para encontrar una fracción equivalente, parte cada pieza en 2 o en 3 partes iguales. Cambia el número de partes y el número de partes coloreadas, pero la cantidad coloreada sigue igual.

Si divides cada parte en 2, el numerador y el denominador se multiplican por 2:
\[
\frac{1}{2}=\frac{1\times2}{2\times2}=\frac{2}{4}
\]
\end{wbidea}

\begin{wbexample}[De $\frac{1}{2}$ a $\frac{2}{4}$]
\textbf{Paso 1.} Dibuja una tira y divídela en 2 partes iguales. Colorea 1 parte: $\frac{1}{2}$.

\textbf{Paso 2.} Divide cada mitad en 2 partes iguales. Ahora hay 4 partes.

\textbf{Paso 3.} Las partes pequeñas coloreadas ahora son 2 de 4: $\frac{2}{4}$.

\textbf{Paso 4.} Escribe la igualdad:
\[
\frac{1}{2}=\frac{1\times2}{2\times2}=\frac{2}{4}
\]

\begin{center}
\wbfracbar[6cm]{1}{2}\\[2pt]
\wbfracbar[6cm]{2}{4}
\end{center}

\textbf{Respuesta:} $\frac{1}{2}=\frac{2}{4}$.
\end{wbexample}

\begin{wbexample}[De $\frac{1}{3}$ a $\frac{2}{6}$]
\textbf{Paso 1.} Dibuja una tira y divídela en 3 partes iguales. Colorea 1 parte: $\frac{1}{3}$.

\textbf{Paso 2.} Divide cada tercio en 2 partes iguales. Ahora hay 6 partes.

\textbf{Paso 3.} La parte coloreada ahora son 2 de 6: $\frac{2}{6}$.

\textbf{Paso 4.} Escribe la igualdad:
\[
\frac{1}{3}=\frac{1\times2}{3\times2}=\frac{2}{6}
\]

\begin{center}
\wbfracbar[6cm]{1}{3}\\[2pt]
\wbfracbar[6cm]{2}{6}
\end{center}

\textbf{Respuesta:} $\frac{1}{3}=\frac{2}{6}$.
\end{wbexample}

\begin{wbexample}[De $\frac{1}{4}$ a $\frac{2}{8}$]
\textbf{Paso 1.} Dibuja una tira y divídela en 4 partes iguales. Colorea 1 parte: $\frac{1}{4}$.

\textbf{Paso 2.} Divide cada cuarto en 2 partes iguales. Ahora hay 8 partes.

\textbf{Paso 3.} La parte coloreada ahora son 2 de 8: $\frac{2}{8}$.

\textbf{Paso 4.} Escribe la igualdad:
\[
\frac{1}{4}=\frac{1\times2}{4\times2}=\frac{2}{8}
\]

\begin{center}
\wbfracbar[6cm]{1}{4}\\[2pt]
\wbfracbar[6cm]{2}{8}
\end{center}

\textbf{Respuesta:} $\frac{1}{4}=\frac{2}{8}$.
\end{wbexample}

\begin{wbtip}[Error común]
No dividas solo la parte sombreada. Debes dividir todas las partes de la tira o del círculo. Si divides cada parte en 2, todas las partes se vuelven más pequeñas.
\end{wbtip}

\begin{wbpractice}
\begin{enumerate}
\item Divide cada parte de $\frac{1}{2}$ en 2. Escribe la fracción equivalente. \wbfracbox \wblines{1}

\begin{center}
\wbfracbar[6cm]{1}{2}
\end{center}

\item Divide cada parte de $\frac{1}{4}$ en 2. Escribe la fracción equivalente. \wbfracbox \wblines{1}

\begin{center}
\wbfracbar[6cm]{1}{4}
\end{center}

\item Divide cada parte de $\frac{1}{3}$ en 2. Escribe la fracción equivalente. \wbfracbox \wblines{1}

\begin{center}
\wbfracbar[6cm]{1}{3}
\end{center}

\item Divide cada parte de $\frac{1}{2}$ en 3. Escribe la fracción equivalente. \wbfracbox \wblines{1}

\begin{center}
\wbfracbar[6cm]{1}{2}
\end{center}

\item Completa: $\frac{2}{4}=\frac{\wbblank[1cm]}{8}$. \wblines{1}

\item Colorea $\frac{2}{4}$ en la primera tira y $\frac{4}{8}$ en la segunda. ¿Valen lo mismo? \wblines{2}

\begin{center}
\wbfracbar[6cm]{0}{4}\\[2pt]
\wbfracbar[6cm]{0}{8}
\end{center}

\item Una barra de chocolate se parte en 3 partes iguales y tomas 1. Después, cada parte se parte en 2. ¿Cuántas piezas pequeñas tomas de 6? \wblines{2}

\item Explica cómo obtienes $\frac{2}{6}$ a partir de $\frac{1}{3}$. \wblines{2}
\end{enumerate}
\end{wbpractice}

\section{Equivalentes en la recta numérica}

Recuerda de la unidad anterior: en la recta numérica divides el espacio de 0 a 1 en partes iguales. Las fracciones equivalentes caen en el mismo punto.

\begin{center}
\wbnumberline[8cm]{2}{1}{1}\\[2pt]
\wbnumberline[8cm]{4}{1}{2}
\end{center}

El punto de $\frac{1}{2}$ y el punto de $\frac{2}{4}$ están en el mismo lugar:
\[
\frac{1}{2}=\frac{2}{4}
\]

También puedes ver $\frac{1}{3}$ y $\frac{2}{6}$.

\begin{center}
\wbnumberline[8cm]{3}{1}{1}\\[2pt]
\wbnumberline[8cm]{6}{1}{2}
\end{center}

\[
\frac{1}{3}=\frac{2}{6}
\]

\begin{wbidea}[El mismo punto]
En la recta numérica, las fracciones equivalentes señalan el mismo punto. Por ejemplo, $\frac{1}{2}$, $\frac{2}{4}$, $\frac{3}{6}$ y $\frac{4}{8}$ están en el mismo lugar.
\end{wbidea}

\begin{wbexample}[$\frac{1}{2}$ y $\frac{2}{4}$ en la recta]
\textbf{Paso 1.} Dibuja una recta de 0 a 1 dividida en 2 partes iguales. Marca $\frac{1}{2}$.

\textbf{Paso 2.} Dibuja otra recta igual de 0 a 1 dividida en 4 partes iguales. Marca $\frac{2}{4}$.

\textbf{Paso 3.} Compara los dos puntos. Caen en el mismo lugar.

\begin{center}
\wbnumberline[8cm]{2}{1}{1}\\[2pt]
\wbnumberline[8cm]{4}{1}{2}
\end{center}

\textbf{Respuesta:} $\frac{1}{2}=\frac{2}{4}$.
\end{wbexample}

\begin{wbexample}[$\frac{1}{3}$ y $\frac{2}{6}$ en la recta]
\textbf{Paso 1.} Dibuja una recta de 0 a 1 dividida en 3 partes iguales. Marca $\frac{1}{3}$.

\textbf{Paso 2.} Dibuja otra recta igual de 0 a 1 dividida en 6 partes iguales. Marca $\frac{2}{6}$.

\textbf{Paso 3.} Observa que los puntos coinciden.

\begin{center}
\wbnumberline[8cm]{3}{1}{1}\\[2pt]
\wbnumberline[8cm]{6}{1}{2}
\end{center}

\textbf{Respuesta:} $\frac{1}{3}=\frac{2}{6}$.
\end{wbexample}

\begin{wbexample}[$\frac{1}{2}$ y $\frac{4}{8}$ en la recta]
\textbf{Paso 1.} Dibuja una recta de 0 a 1 dividida en 2 partes iguales. Marca $\frac{1}{2}$.

\textbf{Paso 2.} Dibuja otra recta igual de 0 a 1 dividida en 8 partes iguales. Marca $\frac{4}{8}$.

\textbf{Paso 3.} Los dos puntos están en el mismo lugar.

\begin{center}
\wbnumberline[8cm]{2}{1}{1}\\[2pt]
\wbnumberline[8cm]{8}{1}{4}
\end{center}

\textbf{Respuesta:} $\frac{1}{2}=\frac{4}{8}$.
\end{wbexample}

\begin{wbtip}[Cuidado]
Para comparar en la recta numérica, las dos rectas deben empezar en 0 y medir lo mismo. Si una recta es más larga, no puedes decir que los puntos están en el mismo lugar.
\end{wbtip}

\begin{wbpractice}
\begin{enumerate}
\item Observa las rectas. ¿$\frac{1}{2}$ y $\frac{3}{6}$ están en el mismo punto? \wblines{1}

\begin{center}
\wbnumberline[7cm]{2}{1}{1}\\[2pt]
\wbnumberline[7cm]{6}{1}{3}
\end{center}

\item Completa: $\frac{1}{2}=\frac{\wbblank[1cm]}{8}$. \wblines{1}

\item En la recta de sextos, marca el punto de $\frac{1}{3}$. \wblines{1}

\wbnumberline[7cm]{6}{1}{}

\item ¿Qué fracción está en el mismo punto que $\frac{2}{4}$: $\frac{1}{2}$, $\frac{1}{4}$ o $\frac{3}{4}$? \wblines{1}

\item Compara: $\frac{3}{6}$ \wbblank[1.5cm] $\frac{1}{2}$. \wblines{1}

\item Sofía camina $\frac{1}{2}$ de kilómetro y su hermano camina $\frac{3}{6}$ de kilómetro por la misma ruta. ¿Quién llegó más lejos? Explica. \wblines{3}

\item En la recta de octavos, marca $\frac{4}{8}$. ¿Qué fracción con denominador 2 está en el mismo punto? \wblines{2}

\wbnumberline[7cm]{8}{1}{}

\item Explica por qué $\frac{2}{6}$ y $\frac{1}{3}$ están en el mismo punto de la recta. \wblines{2}
\end{enumerate}
\end{wbpractice}

\begin{wbreview}
\begin{enumerate}
\item Observa las tiras. ¿$\frac{1}{2}$ y $\frac{2}{4}$ son equivalentes? \wblines{1}

\begin{center}
\wbfracbar[5cm]{1}{2}\\[2pt]
\wbfracbar[5cm]{2}{4}
\end{center}

\item Colorea $\frac{1}{3}$ en el primer círculo y $\frac{2}{6}$ en el segundo. ¿Muestran la misma cantidad? \wblines{2}

\begin{center}
\wbfraccircle[1.1cm]{0}{3}\hspace{0.8cm}\wbfraccircle[1.1cm]{0}{6}
\end{center}

\item Completa: $\frac{1}{2}=\frac{\wbblank[1cm]}{4}$. \wblines{1}

\item Completa: $\frac{1}{4}=\frac{2}{\wbblank[1cm]}$. \wblines{1}

\item Divide cada parte de $\frac{1}{2}$ en 3. Escribe la fracción equivalente. \wbfracbox \wblines{1}

\item Observa las rectas. ¿$\frac{3}{6}$ y $\frac{1}{2}$ están en el mismo punto? \wblines{1}

\begin{center}
\wbnumberline[7cm]{6}{1}{3}\\[2pt]
\wbnumberline[7cm]{2}{1}{1}
\end{center}

\item Completa: $\frac{2}{4}=\frac{\wbblank[1cm]}{8}$. \wblines{1}

\item En una fiesta, Kim comió $\frac{2}{4}$ de un pastel y Alex comió $\frac{1}{2}$ de otro pastel del mismo tamaño. ¿Quién comió más? Explica. \wblines{3}

\item Explica qué significa que dos fracciones sean equivalentes. \wblines{3}
\end{enumerate}
\end{wbreview}
%%% END UNIT
%%% BEGIN ANSWERS
\wbanswers{\wbLPractice\ 4.1}
\begin{enumerate}
\item Sí, $\frac{1}{2}$ y $\frac{2}{4}$ son equivalentes.
\item Se colorea 1 de 3 y 2 de 6. Sí, cubren la misma cantidad.
\item $2$
\item $2$
\item $=$
\item Ninguno pintó más. Ambos pintaron lo mismo: $\frac{1}{2}=\frac{2}{4}$.
\item Respuesta correcta: representan la misma cantidad de un entero del mismo tamaño, aunque tengan diferente número de partes.
\end{enumerate}

\wbanswers{\wbLPractice\ 4.2}
\begin{enumerate}
\item $\frac{2}{4}$
\item $\frac{2}{8}$
\item $\frac{2}{6}$
\item $\frac{3}{6}$
\item $4$
\item Se colorea 2 de 4 y 4 de 8. Sí, valen lo mismo.
\item Tomas 2 piezas de 6: $\frac{1}{3}=\frac{2}{6}$.
\item Se divide cada tercio en 2 partes iguales; hay 6 partes en total y 2 quedan coloreadas.
\end{enumerate}

\wbanswers{\wbLPractice\ 4.3}
\begin{enumerate}
\item Sí, están en el mismo punto.
\item $4$
\item Se marca $\frac{2}{6}$.
\item $\frac{1}{2}$
\item $=$
\item Ninguno llegó más lejos. Los dos caminaron lo mismo: $\frac{1}{2}=\frac{3}{6}$.
\item Se marca $\frac{4}{8}$, que está en el mismo punto que $\frac{1}{2}$.
\item Cada tercio se parte en 2 sextos, y 2 sextos ocupan el mismo espacio que 1 tercio.
\end{enumerate}

\wbanswers{\wbLReview}
\begin{enumerate}
\item Sí, son equivalentes.
\item Se colorea 1 de 3 y 2 de 6. Sí, muestran la misma cantidad.
\item $2$
\item $8$
\item $\frac{3}{6}$
\item Sí, están en el mismo punto.
\item $4$
\item Ninguno comió más. Ambos comieron lo mismo: $\frac{2}{4}=\frac{1}{2}$.
\item Respuesta correcta: son fracciones que representan la misma cantidad de un entero del mismo tamaño.
\end{enumerate}
%%% END ANSWERS